\begin{figure}[!htbp]
\centering
\begin{tikzpicture}[beach diagram]
  \node[beach heading] at (3.4,6.1) {(a) ray の最初の命中と遮蔽};
  \draw[thick] (0.2,5.1)--node[above] {照射開口 $A_{\mathrm{ap}}$} (6.5,5.1);
  \fill[black!8] (0,0.8)--(0,1.3)--(2.8,1.3)--(2.8,3.2)
    --(4.1,3.2)--(4.1,1.3)--(6.7,1.3)--(6.7,0.8)--cycle;
  \draw (0,1.3)--(2.8,1.3)--(2.8,3.2)--(4.1,3.2)--(4.1,1.3)--(6.7,1.3);
  \fill[black!20] (1.2,1.3)--(2.8,3.2)--(2.8,1.3)--cycle;
  \foreach \x in {2.7} {
    \draw[beach arrow,BeachOrange] (\x,5.1)--(\x-2.45,1.3);
  }
  \draw[beach arrow,BeachOrange] (4.1,5.1)--(2.875,3.2);
  \draw[beach arrow,BeachOrange] (5.3,5.1)--(4.075,3.2);
  \draw[beach arrow,BeachOrange] (6.5,5.1)--(4.1,1.378);
  \node[beach note,fill=white,inner sep=1pt] at (1.65,1.85) {影};
  \node[beach dot] at (2.875,3.2) {};
  \draw[beach arrow,BeachBlue] (2.9,3.3) .. controls (3.1,4.05) and (3.6,4.05) .. (3.8,4.3);
  \node[anchor=west] at (3.85,4.25) {光電子};
  \node[beach note,text width=65mm] at (3.35,0)
    {$A_{\mathrm{proj}}=A_{\mathrm{ap}}|\widehat{\vect d}\cdot\vect n_{\mathrm{in}}|$};
  \node[beach heading] at (11.3,6.1) {(b) 電子一粒子の電荷収支};
  \node[beach box,text width=49mm,fill=BeachOrange!6] (emit) at (11.3,4.8)
    {放出元 $i$：$-qw=+|q|w$};
  \node[beach box,text width=49mm,fill=BeachBlue!6] (hit) at (11.3,2.8)
    {帰還先 $j$：$+qw=-|q|w$};
  \draw[beach arrow,BeachBlue] (emit)--node[right] {再吸収} (hit);
  \node[beach note,text width=58mm] at (11.3,1.35)
    {同じ要素へ戻る：相殺\\別の要素へ戻る：電荷移送};
  \node[beach note,text width=58mm] at (11.3,0)
    {escape：放出反作用が残る\\箱外への電子電荷は表面へ deposit しない};
\end{tikzpicture}
\caption[照射 ray と光電子の電荷収支]{照射 ray と光電子の電荷収支。左は平行な ray が凹凸に遮られる断面模式図で、
命中しない ray は粒子を生成しない。右は $q<0$ の電子についての符号を示す。
放出後の軌道が再吸収と escape を決め、放出重みに事前の escape 率は掛けない。}
\label{fig:photoemission}
\end{figure}
